\subsection{Inclusion-Exclusion, Pigeonhole, and Combinatorial Proof}

\begin{definition} \label{def:multisets-and-multiplicity} \hlt{Multisets}\\
Let $S$ be a set. A \hlt{multiset on $S$} is a pair $(S_0,m)$ such that
\begin{enumerate}[label=(\roman*)]
\setlength{\itemsep}{0pt}
\item $S_0\subseteq S$ is the \hlt{underlying set};
\item $m:S_0\to\Z_{>0}$ is the \hlt{multiplicity function};
\item $m(s)$ is the \hlt{multiplicity} of $s\in S_0$;
\item if $S_0$ is finite, the \hlt{size} is $\lvert m\rvert\defeq\sum_{s\in S_0}m(s)$.
\end{enumerate}
Equivalently, record $(S_0,m)$ as $\{(s,m(s)):s\in S_0\}$.
\end{definition}

\begin{example} \label{ex:multiset-eigenvalue-bookkeeping} \hlt{Multiset Bookkeeping}\\
The list $\{a,a,b,b,b,c\}$ is the multiset on $\{a,b,c\}$ with multiplicities $2,3,1$.
This language is useful whenever order is irrelevant but repetition matters, such as eigenvalues counted with algebraic multiplicity.
\end{example}

\hlt{Placeholder}: Add inclusion-exclusion, pigeonhole arguments, and combinatorial proof as equality by double counting.
